\chapter{The Marseille Minor Extension}

The Mage's Guild's extended system is English cartomancy with a Marseille Minor
extension. The English playing-card system is the normative base. Marseille
Minors may add compatible detail; they never override English meanings.

\section{The precedence rule}

English suits remain the domain authority, English numerology remains the rank
authority, and English courts remain the interpersonal authority. If a Marseille
Minor reading conflicts with any of these, the English system replaces the
conflicting interpretation. The Marseille observation may be recorded as a
native observation, but it does not enter the canonical blended reading.

The Marseille Majors are excluded from this extension. They belong to the
Christmas Tree programming substrate. Golden Dawn elemental correspondences,
Qabalistic attributions, and Smith--Waite meanings are also excluded here.

\section{Suit bridges}

\begin{longtable}{p{0.22\textwidth}p{0.22\textwidth}p{0.46\textwidth}}
\toprule
English & Marseille & Compatible extension \\
\midrule
Hearts & Cups & feeling, affection, home, intimacy, emotional flow \\
Diamonds & Coins & material value, resources, exchange, property, practical embodiment \\
Clubs & Batons & effort, enterprise, growth, initiative, directed activity \\
Spades & Swords & challenge, separation, consequence, truth, boundaries, conflict \\
\bottomrule
\end{longtable}

These are bridges, not replacements. Diamonds also govern communication and
news in the English system; Coins do not erase that meaning. Swords can describe
clarity and decision as well as difficulty; they are not automatically negative.

\section{Pip geometry}

The Marseille pip card contributes visual structure: symmetry, axis, enclosure,
crossing, density, repetition, opening, or closure. The English number still
determines the operation. Marseille geometry determines how that operation is
expressed in material or relational form.

The Guild's standing-wave, shell-shape, and baseline-physics readings are a
developing interpretive layer. They are not silently presented as historical
Marseille doctrine.

\section{The four courts}

Marseille adds the Cavalier or Knight to the English Jack, Queen, and King.
The compatible extension is:

\begin{itemize}
\item Valet/Jack: youth, messenger, developing agency, initial movement;
\item Cavalier/Knight: a force in transit, mission, active transmission;
\item Queen: mature relational or inward expression;
\item King: mature authority, stewardship, expertise, or external power.
\end{itemize}

The Knight is an added Marseille role, not a revision of English court meanings.

\section{Cross-reading}

An English card may be expanded by a Marseille Minor, but the parent English
meaning remains intact. A Marseille Minor may be grounded by an English card,
which supplies the worldly domain. Every expansion records its parent, deck,
position, order, and interpretive priority.

\section{The Minor as a second skin}

The Marseille Minor does not give the English card a new soul. It gives the
English operation a second skin: a visible arrangement of vessels, blades,
branches, or coins through which the worldly event may show its body.

The English card asks what is happening. The Marseille pip may help show how the
event is holding itself together. Is the force enclosed or open? Does it cross
itself? Does it gather at an axis, repeat around a center, leave a channel, or
press against a boundary? These observations are useful when they clarify the
English operation and remain subordinate when they do not.

\section{A--Ten: operation and embodiment}

The following bridge is the working heart of the extension. The English number
is the operation. The Marseille pip is the local body of that operation.

\begin{longtable}{p{0.12\textwidth}p{0.32\textwidth}p{0.46\textwidth}}
\toprule
English rank & English operation & Compatible Marseille question \\
\midrule
Ace & maximum force, concentration & Where is the suit gathered into one vessel, axis, or undivided presence? \\
Two & relation, encounter & What is meeting, mirroring, crossing, or being held apart? \\
Three & multiplication, first social form & What third element, branch, support, or complication is being made? \\
Four & consolidation, stability & What has become enclosed, squared, sheltered, or stuck? \\
Five & disruption, turning & Where does the established pattern break, open, or change direction? \\
Six & gradual resolution, repair & What channel, exchange, bridge, or rebalancing carries the pattern onward? \\
Seven & test, obstacle, careful handling & What crossing, obstruction, asymmetry, or pressure demands skill? \\
Eight & movement, approach, arrival & What line, repetition, current, or force is travelling toward completion? \\
Nine & ripening, intensity near completion & Where has the pattern become dense, abundant, strained, or almost full? \\
Ten & completion, full result & What has become whole enough to harvest, close, transform, or set down? \\
\bottomrule
\end{longtable}

This table is not a secret Marseille dictionary. It is a set of questions for
the eye. If the English Five of Diamonds describes a disruptive message, the
Coins or pip geometry may show whether that message is an offer opened, a
resource divided, a contract crossed, or a material arrangement losing its
symmetry. The English meaning remains the sentence. The pip shows the weather
inside the sentence.

\section{Pip geometry as standing wave}

The reader may attend to the pips as a small field of forces:

\begin{itemize}
\item an axis may show a line of decision, transmission, or balance;
\item a mirror or symmetry may show correspondence, partnership, or a pattern
      repeating across a divide;
\item an enclosure may show protection, containment, privacy, or stagnation;
\item a crossing may show conflict, exchange, interruption, or a necessary
      meeting;
\item density may show concentration, burden, abundance, or pressure;
\item an opening may show access, release, invitation, or exposure;
\item a broken symmetry may show imbalance, asymmetrical labor, or a new form
      trying to emerge;
\item a center may show the place where the English operation is being held.
\end{itemize}

These are not automatic meanings. The reader first names the English operation,
then notices which geometry makes that operation more particular. A symmetrical
Two of Cups does not erase the English Two's possibility of opposition. It asks
whether the relation is balanced, mirrored, or merely facing itself. A crossed
Five of Swords does not make disruption noble or wicked. It shows where the
turn has become a crossing that must be handled.

The Guild calls this a standing-wave reading because the pip field may reveal a
pattern that persists while the English operation moves through it. The name is
a Guild metaphor, not a claim about historical Marseille terminology or literal
physics.

\section{Material bodies of the four bridges}

\subsection{Cups: the vessel}

Cups give Hearts a vessel-language. They show what feeling can hold, pour,
receive, overflow, or leave empty. Use them to ask whether affection has a home,
whether grief has somewhere to go, or whether a relationship is offering a cup
that the English card says must still be tested.

\subsection{Coins: the weight}

Coins give Diamonds a weight-language. They show what can be counted, exchanged,
owned, spent, carried, or made durable. Coins can make a message feel like a
document, an offer like a price, or a practical question like a thing that must
be placed on a table and handled.

\subsection{Batons: the reach}

Batons give Clubs a reach-language. They show extension, leverage, growth,
direction, effort, and the work of joining one point to another. A Baton can be
the tool, branch, staff, or living shoot through which enterprise finds a path.

\subsection{Swords: the edge}

Swords give Spades an edge-language. They show separation, clarity, decision,
conflict, boundaries, and the cost of cutting one thing away from another. A
Sword is not automatically cruelty. It may be the instrument that makes an
honest shape possible.

\section{The Knight enters the room}

The Marseille Knight or Cavalier adds a distinct motion to the English courts.
The Jack is developing agency or a messenger at the beginning of movement. The
Knight is agency in transit: a mission, rider, transmission, delivery, journey,
or force carrying one condition into another. The Queen gives the suit mature
relational or inward intelligence. The King gives it mature authority,
stewardship, expertise, or external power.

The Knight is especially useful in recursive readings. If an English court is
the parent and a Marseille Knight appears as child, ask what that court is
carrying, where it is going, and whether the message survives the journey. If
the English card is a Jack, the Knight may show the next maturation of its
agency. If the English card is a King or Queen, the Knight may show delegated
power: authority in motion rather than authority seated.

\section{Precedence in practice}

The conflict rule becomes clearest through examples:

\begin{enumerate}
\item An English Five of Diamonds beside a Marseille Five of Coins shows a
      material or communicative disruption. If the pip arrangement appears
      harmonious, record that as a possible resource within the disruption; do
      not erase the English turn.
\item An English Seven of Spades expanded by a Seven of Swords may show a
      difficult boundary or test made visible through crossing pips. If a reader
      prefers to call it a triumph, that is not yet licensed by the English
      Seven; the card still demands careful handling.
\item An English Two of Hearts expanded by a Two of Cups may show a relation
      with emotional capacity and mirrored feeling. If the Cups appear divided,
      the English Two still permits separation or opposition. The geometry
      sharpens the question; it does not choose the answer alone.
\item An English Eight of Clubs expanded by an Eight of Batons may show effort
      travelling toward a goal. A broken Baton line can add delay or diversion,
      but cannot remove the English Eight's movement without support from
      position and neighboring cards.
\end{enumerate}

In every case, write the English observation first. Add the Marseille
observation second. If the two cannot live together, preserve both in the raw
record and let the English reading remain canonical.

\section{A guide may speak through the pip field}

A guide may draw the reader's eye to a crossing, an empty space, a repeated
axis, or a pip that seems to become a door. The reader should not force the
geometry to speak merely because a guide was invited. Ask instead whether the
same kind of emphasis has become recognizable across time. The guide may say,
``look here,'' while the English number and suit say what is being looked at.

Some readers will experience this as an old hand resting on the spread: not a
voice dictating a result, but a patient pressure toward the detail that matters.
Let that hand remain a companion to seeing. The pip field is still observed,
the English authority is still honored, and the reader remains free to close the
book, set down the card, or decline the invitation.
